% Part: proof-theory
% Chapter: natural-deduction
% Section: sequents

\documentclass[../../../include/open-logic-section]{subfiles}

\begin{document}

\olfileid{pt}{nat}{seq}

\olsection{સિક્વન્ટો સાથેનું સ્વાભાવિક નિગમન: \Log{N2c} અને~\Log{N2i}}

\Log{N1c} અથવા~\Log{N1i}માંની
!!^a{proof}~$\delta$
બતાવે છે કે તેનું
અંતિમ-!!{formula}~$!A$
તેની ખુલ્લી ધારણાઓમાંથી
ફલિત થાય છે.
\sourcecorrection{OLMD-237}{મૂળમાં અહીં G1c/G1i લખાયું છે, પણ ખુલ્લી ધારણાવાળી સૂત્ર-સાબિતીઓની ચર્ચા N1c/N1i અંગે છે; G1 તંત્ર સિક્વન્ટ-કલનનાં છે.}
જો~$\Gamma$ એવા
!!{formula}s~$!B$નો
ગણ હોય કે~$!B^x$
એ~$\delta$માં
ખુલ્લી ધારણા હોય,
તો આને
$\delta \Proves \Gamma
\Sequent !A$ એમ
લખી શકાય.
આથી સિક્વન્ટો માટે
પણ સ્વાભાવિક નિગમનનું
તંત્ર ઘડી શકાય એવું
સૂઝે છે.
આવા તંત્રમાં દરેક
સિક્વન્ટ જણાવે છે કે
જમણી બાજુનું !!{formula}
કઈ ખુલ્લી ધારણાઓ
પર આધારિત છે,
એટલે એ સિક્વન્ટમાં
પૂરી થતી ઉપ-!!{proof}માં
કઈ ધારણાઓ ખુલ્લી છે.
આ સિક્વન્ટો સાથેનું,
અથવા ``સિક્વન્ટ-શૈલીનું'',
સ્વાભાવિક નિગમન છે;
મળતાં તંત્રોને
\Log{N2c} અને~\Log{N2i}
કહીએ છીએ.

\Log{N1} અને~\Log{N2}
વચ્ચેનો અનુરૂપ સંબંધ
વધુ સીધો બનાવવા
ડાબી બાજુના ગણો~$\Gamma$
ચિહ્નિત
!!{formula}s~$x:!B$
ધરાવે એમ લઈએ.
\Log{N1}ના નિયમો
આધારવાક્યોનાં
!!{formula}s પર,
એટલે તાત્કાલિક
ઉપ-!!{proof}sના
અંતિમ-!!{formula}s પર,
જ કામ કરે છે;
તેથી~\Log{N2}ના
નિયમો સિક્વન્ટના
ઉત્તરાંગ પર જ
કામ કરે છે.
આથી સિક્વન્ટની
ડાબી બાજુનાં
ચિહ્નિત સૂત્રોને
ફક્ત \emph{સંદર્ભ}
કહી શકાય.

\Log{N2}ની સાબિતીઓમાં
ધારણાઓને બદલે
સૌથી ઉપર
આ સ્વરૂપનાં
સ્વયંસિદ્ધ સિક્વન્ટો છે:
\[x:!A \Sequent !A.\]
આ નિયમો~\Log{N1}ના
નિયમોને ચોક્કસ અનુરૂપ
છે અને
\olref{tab:N2}માં
આપેલા છે.

\subfile{rules-N2}

\Log{N1c} અને~\Log{N1i}ની
જેમ અહીં પણ
\Intro{\lif}, \Intro{\lnot},
\Elim{\lor} અને
\Elim{\lexists} નિયમો
ધારણાઓ મુક્ત કરી
શકે, પણ એમ કરવું
જરૂરી ન હોય,
એવી છૂટ જોઈએ.
તેથી~\Log{N2c}
અને~\Log{N2i}માં
સંબંધિત આધારવાક્યના
સંદર્ભમાં અનુરૂપ
!!{formula}s~$!A$,
$!B$ અને~$!A(a)$
હાજર ન હોય તોપણ
એ નિયમો યોગ્ય રીતે
લાગુ પાડી શકાય છે.

\begin{prop}
જો~$\delta$ એ~$!A$ની
\Log{N1c} અથવા~\Log{N1i}માંની
!!a{proof} હોય,
તો~\Log{N2c}
(\Log{N2i})માં
$\Gamma \Sequent !A$ની
!!a{proof}~$\delta'$
છે. અહીં~$\Gamma$
એ એવા બધા~$x:!B$નો
ગણ છે કે~$!B^x$
એ~$\delta$ની
ખુલ્લી ધારણા છે.
\end{prop}

\begin{proof}
  $n=\pheight{\delta}$
  પર આગમન કરીએ.
  જો~$n=0$, તો~$\delta$
  એક જ ખુલ્લી
  ધારણા~$!A^x$ છે.
  ત્યારે~$\Gamma = \{x:!A\}$
  અને~$\Gamma
  \Sequent !A$ એ
  \Log{N2c} અથવા~\Log{N2i}
  નું સ્વયંસિદ્ધ છે.

  જો~$n>0$, તો~$\delta$ના
  છેલ્લા અનુમાન પ્રમાણે
  કિસ્સા જુદા પાડીએ.
  અહીં માત્ર એક
  કિસ્સો જોઈએ;
  બાકીના કસરત
  તરીકે રાખીએ છીએ.

  માનો કે છેલ્લું
  અનુમાન~$\Intro{\lif}$
  છે. ત્યારે~$\delta$નો
  અંત આ રીતે થાય છે:
  \[
  \AxiomC{$\Discharge{!B}{x}$}
  \RightLabel{$\delta_1$}
  \DeduceC{$!C$}
  \DischargeRule{\Intro{\lif}}{x}
  \UnaryInfC{$!B \lif !C$}
  \DisplayProof
  \]
  આગમન-પરિકલ્પનાથી
  \Log{N2c} (\Log{N2i})માં
  $\Gamma' \Sequent !C$ની
  !!a{proof}~$\delta_1'$
  છે, જ્યાં~$\Gamma'$
  એ~$\delta_1$માં
  ખુલ્લી ચિહ્નિત
  ધારણાઓનો ગણ છે.
  જો~$!B^x$,
  $\delta_1$ની ખુલ્લી
  ધારણા હોય, તો
  \Intro{\lif} અનુમાન
  તેને મુક્ત કરે છે,
  અને~$\delta$ની
  ખુલ્લી ધારણાઓ
  $\Gamma =
  \Gamma' \setminus \{x:!B\}$
  છે. એટલે~$\delta_1'$
  એ~$x:!B, \Gamma \Sequent !C$ની
  !!a{proof} છે,
  અને~$\delta'$ને
  આ રીતે લઈ શકાય:
  \[
  \AxiomC{}
  \RightLabel{$\delta_1'$}
  \Deduce$x:!B, \Gamma \fCenter !C$
  \RightLabel{\Intro{\lif}}
  \UnaryInf$\Gamma \fCenter !B \lif !C$
  \DisplayProof
  \]
  \sourcecorrection{OLMD-238}{મૂળના આ N2 વૃક્ષમાં પૂર્વાંગમાં $x:B$ લખાયું છે; તરતની દલીલ અને ચિહ્નિત સૂત્રની વ્યાખ્યા પ્રમાણે $x:!B$ જોઈએ.}
  જો~$!B^x$,
  $\delta_1$ની ખુલ્લી
  ધારણા ન હોય,
  તો~\Intro{\lif}
  કોઈ ધારણા મુક્ત
  કરતું નથી, અને
  $\delta$ તથા~$\delta_1$ની
  ખુલ્લી ધારણાઓ સમાન છે.
  \Log{N2c} તથા~\Log{N2i}માં
  \Intro{\lif}ના
  આધારવાક્યના સંદર્ભમાં
  $x:!B$ હાજર હોવું
  જરૂરી ન હોવાથી,
  $\delta'$ આ રીતે
  લઈ શકાય:
  \[
  \AxiomC{}
  \RightLabel{$\delta_1'$}
  \Deduce$\Gamma \fCenter !C$
  \RightLabel{\Intro{\lif}}
  \UnaryInf$\Gamma \fCenter !B \lif !C$
  \DisplayProof
  \]

  $\delta$નો અંત બીજા
  કોઈ નિયમમાં થાય
  તે કિસ્સા સમાન છે.
\end{proof}

\begin{prop}
જો~$\delta$ એ
\Log{N2c} અથવા~\Log{N2i}માં
$\Gamma \Sequent !A$ની
!!a{proof} હોય,
તો~\Log{N1c} (\Log{N1i})
માં એવી
!!a{proof}~$\delta'$
છે જેનું
અંતિમ-!!{formula}~$!A$
છે અને
દરેક
$x:!B \in \Gamma$
માટે~$!B^x$ તેની
ખુલ્લી ધારણા છે.
\end{prop}

\begin{proof}
  ફરી~$\delta$ની
  !!{height}~$n$ પર
  આગમન કરીએ.
  જો~$n=0$,
  તો~$\delta$ એ
  સ્વયંસિદ્ધ
  $x:!A \Sequent !A$ છે.
  !!{proof}~$\delta'$
  એ ધારણા~$!A^x$ છે.

  જો~$n>0$, તો~$\delta$ના
  છેલ્લા અનુમાન પ્રમાણે
  કિસ્સા જુદા પાડીએ.
  દાખલા તરીકે, અનુમાન
  $\Intro{\lif}$ હોય,
  તો~$\delta$નો અંત
  આ બેમાંથી એક રીતે થાય છે:
  \[
  \bottomAlignProof
  \AxiomC{}
  \RightLabel{$\delta_1$}
  \Deduce$x:!B, \Gamma \fCenter !C$
  \RightLabel{\Intro{\lif}}
  \UnaryInf$\Gamma \fCenter !B \lif !C$
  \DisplayProof
  \quad\text{ અથવા }\quad
  \bottomAlignProof
  \AxiomC{}
  \RightLabel{$\delta_1$}
  \Deduce$\Gamma \fCenter !C$
  \RightLabel{\Intro{\lif}}
  \UnaryInf$\Gamma \fCenter !B \lif !C$
  \DisplayProof
  \]
  આગમન-પરિકલ્પનાથી
  \Log{N1c} (\Log{N1i})
  માં~$!C$ની
  !!a{proof}~$\delta_1'$
  છે. પહેલા કિસ્સામાં
  $!B^x$ એ તેની
  ખુલ્લી ધારણા છે;
  બીજામાં નથી.
  !!{proof}~$\delta'$ને
  આ રીતે લઈ શકાય:
  \[
  \bottomAlignProof
  \AxiomC{$\Discharge{!B}{x}$}
  \RightLabel{$\delta_1'$}
  \DeduceC{$!C$}
  \DischargeRule{\Intro{\lif}}{x}
  \UnaryInfC{$!B \lif !C$}
  \DisplayProof
  \quad\text{ અથવા }\quad
  \bottomAlignProof
  \AxiomC{}
  \RightLabel{$\delta_1'$}
  \DeduceC{$!C$}
  \RightLabel{\Intro{\lif}}
  \UnaryInfC{$!B \lif !C$}
  \DisplayProof
  \]
  અનુક્રમે.
  \sourcecorrection{OLMD-239}{મૂળના અંતિમ બે N1 વૃક્ષોમાં $\delta_1$ લખાયું છે, જે N2 ઉપસાબિતી છે; આગમનથી મળતી N1 સાબિતી $\delta_1'$ છે. પહેલીમાં $!B^x$ ખુલ્લું છે તે પણ એ જ ઉપસાબિતી માટે કહેવાય છે.}
\end{proof}

\end{document}
